% === Chapter 5: Thread 5 — Multimodal Post-Training ===
\section{Thread 5: Multimodal Post-Training}
\label{sec:t5}

% ============================================================
% 5.1 Problem Origin
% ============================================================
\subsection{Problem Origin: From Text-Only to Vision-Language}
\label{sec:t5-origin}

The core dilemma facing pre-trained vision-language models (VLMs) is highly similar to, yet more acute than, that of text-only LLMs:
models trained solely through image-text contrastive learning (e.g., CLIP \cite{radford2021clip}) or
image-text matching pre-training (e.g., BLIP-2 \cite{li2023blip2}),
while possessing a foundational capability for vision-language alignment,
can neither follow open-ended multimodal instructions nor generate structured long-form responses.
In other words, the pre-training stage endows VLMs with the ability to \textbf{perceive},
while the post-training stage must endow them with the ability to \textbf{reason, follow instructions, and stay faithful}.

\begin{problembox}
How can the mature SFT, RLHF/DPO, and other techniques developed in text-only post-training
be transferred to the vision-language joint modeling setting?
This transfer faces three major challenges:
\begin{enumerate}[leftmargin=2em]
  \item \textbf{Modality Alignment}:
    Visual features and text tokens reside in different embedding spaces.
    How should projection or fusion mechanisms be designed to enable the LLM to ``understand'' visual information?
  \item \textbf{Data Construction}:
    The annotation cost of multimodal instruction-following data is far higher than that of pure text.
    How can existing vision-language datasets be leveraged, or strong models be used to automatically generate high-quality training data?
  \item \textbf{Multimodal Hallucination}:
    VLMs generate descriptions inconsistent with image content (e.g., claiming the presence of objects that do not exist in the image).
    While this problem exists in text-only LLMs, it is significantly more pronounced and easier to measure in multimodal settings.
\end{enumerate}
\end{problembox}

\noindent
From the perspective of technical evolution, multimodal post-training directly inherits and extends
the core paradigm of instruction tuning from Thread~1
\crossref{1}{sec:t1},
while also introducing the methodology of preference optimization from Thread~2
\crossref{2}{sec:t2}
into the vision-language setting to address hallucination and alignment problems.
The unique contribution of this thread is that
it not only addresses alignment challenges in multimodal settings,
but also provides new insights for text-only post-training in return ---
for example, research on multimodal hallucination has driven the development of factual grounding techniques,
and training methods for visual grounding have inspired citation generation in text-only LLMs.

% ============================================================
% 5.2 Foundation Methods
% ============================================================
\subsection{Foundation Methods: From Flamingo to Visual Instruction Tuning}
\label{sec:t5-foundation}

This section provides an in-depth analysis of three foundational works:
Flamingo \cite{alayrac2022flamingo} established the architectural paradigm for VLMs,
LLaVA \cite{liu2023llava} pioneered visual instruction tuning,
and InstructBLIP \cite{dai2023instructblip} systematically advanced instruction-aware visual feature extraction.

% --- Flamingo ---
\subsubsection{Flamingo: Visual Language Model for Few-Shot Learning}
\label{sec:t5-flamingo}

\landmark{alayrac2022flamingo}
Flamingo \cite{alayrac2022flamingo} is a vision-language model proposed by DeepMind in 2022.
It was the first to demonstrate that by injecting visual information into a frozen LLM,
powerful few-shot multimodal capabilities can be achieved without fine-tuning.
Flamingo laid the architectural foundation for all subsequent VLM post-training work.

\paragraph{Architecture.}
Flamingo's architecture consists of three core components:

\begin{enumerate}[leftmargin=2em]
  \item \textbf{Vision Encoder}:
    A pre-trained and frozen NFNet-F6 serves as the vision encoder,
    trained through contrastive learning on image-text pairs,
    encoding input images (or video frames) into spatial features.

  \item \textbf{Perceiver Resampler}:
    This is one of Flamingo's key innovations.
    The vision encoder produces a \textit{variable number} of visual features for images/videos of different resolutions.
    The Perceiver Resampler uses a fixed set of \textit{learned latent queries} (64 in total)
    to perform cross-attention over these features,
    mapping an arbitrary number of visual tokens to a fixed set of 64 visual output tokens.
    The core advantages of this design are:
    (1) it decouples the dependency between the number of visual features and the LLM input length;
    (2) it provides a unified interface for handling multi-image and video inputs.

  \item \textbf{Gated Cross-Attention Dense (XATTN-DENSE) Layers}:
    Before each self-attention layer in the frozen LLM (Chinchilla),
    a new cross-attention layer is inserted, allowing text tokens to attend to visual tokens.
    The key design is \textbf{tanh gating}:
    the output of each cross-attention layer is multiplied by a learnable scalar $\tanh(\alpha)$,
    where $\alpha$ is initialized to 0.
    This ensures that at the beginning of training, the newly inserted visual modules have zero influence on the LLM's output,
    thereby protecting the language capabilities acquired during pre-training from being disrupted.
    During training, the model gradually learns to leverage visual information.
\end{enumerate}

\paragraph{Training.}
Flamingo's training strategy embodies the essence of pragmatic engineering:
both the vision encoder and the LLM remain frozen,
and only the Perceiver Resampler, gated cross-attention layers, and layer norms are trained.
The training data comprises three sources:
(1) M3W (MultiModal MassiveWeb) --- 43M web pages with interleaved image-text content crawled from the internet;
(2) LTIP (Long Text \& Image Pairs) --- 312M image-text pairs;
(3) VTP (Video \& Text Pairs) --- 27M video-text pairs.
The unique value of the M3W data is that it naturally provides \textit{interleaved image-text sequences},
enabling Flamingo to handle multiple images in in-context learning.

\paragraph{Results.}
The 80B-parameter Flamingo established new few-shot state-of-the-art results across 16 multimodal benchmarks.
Even more strikingly, a 32-shot few-shot Flamingo surpassed prior fully fine-tuned SOTA models on 6 benchmarks.
This result powerfully demonstrated the effectiveness of the frozen LLM + visual adapter paradigm.
Flamingo's architectural vision quickly spawned open-source reproductions:
IDEFICS \cite{laurencon2023idefics}, built on the publicly available OBELICS dataset,
provided an open-source variant of Flamingo as a reproducible baseline for subsequent VLM research.
Meanwhile, PaLM-E \cite{driess2023palme} extended the Flamingo-style frozen LLM + visual adapter paradigm
to \textbf{embodied settings} --- unifying robot sensor inputs, scene images, and natural language instructions
into the input space of PaLM 562B,
demonstrating that multimodal LLMs can not only ``understand images'' but can also directly drive robots to execute tasks in the physical world
(\crossref{6}{sec:t6}).

\evolution{Separately trained vision model + language model}
  {Unified VLM via Frozen LLM + Visual Adapters}
  {Perceiver Resampler + Gated XATTN}

% --- LLaVA ---
\subsubsection{LLaVA: Visual Instruction Tuning}
\label{sec:t5-llava}

\landmark{liu2023llava}
If Flamingo established the architectural paradigm for VLMs,
then LLaVA (Large Language and Vision Assistant) \cite{liu2023llava} pioneered
the \textbf{visual instruction tuning} training paradigm ---
that is, by performing SFT on multimodal instruction-following data,
it transformed a VLM from a visual question-answering tool into a general-purpose multimodal AI assistant.
This work directly transferred the core idea of instruction tuning from Thread~1
\crossref{1}{sec:t1}
to the vision-language setting.

\paragraph{Architecture.}
LLaVA's architecture is renowned for its simplicity, consisting of three components:
\begin{itemize}[leftmargin=2em]
  \item \textbf{Vision Encoder}:
    A pre-trained CLIP ViT-L/14 (resolution 224$\times$224),
    using the features from the layer before the last transformer layer as the visual representation.
  \item \textbf{Projection Layer}:
    A trainable linear projection matrix $\mathbf{W}$
    that maps CLIP visual features into the LLM's word embedding space:
    $\mathbf{H}_v = \mathbf{W} \cdot \mathbf{Z}_v$,
    where $\mathbf{Z}_v$ is the visual feature sequence output by CLIP.
    This is the sole bridge connecting the visual and language modalities in the entire architecture.
  \item \textbf{LLM}:
    Vicuna \cite{chiang2023vicuna} (an instruction-tuned model based on LLaMA),
    which receives the concatenated sequence of visual tokens and text tokens as input.
\end{itemize}

\noindent
Compared to Flamingo's Perceiver Resampler + gated cross-attention design,
LLaVA's single linear projection layer appears extremely simple.
Yet it is precisely this ``radical simplicity'' that made LLaVA the most easily reproducible VLM baseline,
spawning an extensive body of follow-up work.

\paragraph{Data Construction via GPT-4.}
One of LLaVA's most influential contributions is its method for automatically generating multimodal instruction-following data using GPT-4.
Specifically, the researchers fed captions and bounding box information from COCO \cite{lin2014coco} images
into GPT-4 in pure text form (note: at that time, GPT-4 had no visual capabilities),
and through carefully designed prompts, had GPT-4 generate three types of data:

\begin{itemize}[leftmargin=2em]
  \item \textbf{Conversation} (58K samples):
    Multi-turn question-answering dialogue data centered around image content.
  \item \textbf{Detailed Description} (23K samples):
    Data providing detailed, comprehensive descriptions of image content.
  \item \textbf{Complex Reasoning} (77K samples):
    Question-answering data requiring deep reasoning about image content.
\end{itemize}

\noindent
This ``text-only GPT-4 + image metadata $\rightarrow$ multimodal training data'' paradigm
completely bypassed the high cost of multimodal data annotation,
laying the methodological foundation for data scaling in subsequent works such as
ShareGPT4V \cite{chen2023sharegpt4v} and LLaVA-1.5 \cite{liu2023llava15}.

\paragraph{Two-Stage Training.}
LLaVA's training is divided into two stages:

\begin{enumerate}[leftmargin=2em]
  \item \textbf{Stage 1 --- Feature Alignment Pre-Training}:
    Both the vision encoder and LLM remain frozen,
    and only the projection matrix $\mathbf{W}$ is trained.
    Using 595K image-text pairs filtered from CC3M,
    the training objective is standard next-token prediction (loss computed only on text tokens).
    The purpose of this stage is to enable the projection layer to learn to map visual features into the semantic space that the LLM can understand.

  \item \textbf{Stage 2 --- End-to-End Visual Instruction Tuning}:
    The vision encoder remains frozen,
    while the projection matrix $\mathbf{W}$ and the LLM are both unfrozen and jointly trained.
    Using 158K multimodal instruction-following data generated by GPT-4,
    the training objective remains next-token prediction, but the input is now a complete multimodal instruction.
    This stage enables the model to learn to follow multimodal instructions and generate structured responses.
\end{enumerate}

\paragraph{Results \& Impact.}
On LLaVA-Bench (In-the-Wild),
LLaVA achieved 85.1\% relative score as evaluated by GPT-4 (text-only + image caption input).
More importantly, LLaVA established an extremely influential VLM training paradigm:
\textit{CLIP encoder + simple projection + LLM + two-stage training}.
The subsequent LLaVA-1.5 \cite{liu2023llava15} replaced the linear projection with an MLP,
introduced higher resolution (336$\times$336), and incorporated academic-task-oriented VQA data,
significantly improving performance across multiple benchmarks.
This continuous evolution of the LLaVA family
(LLaVA $\rightarrow$ LLaVA-1.5 $\rightarrow$ LLaVA-NeXT $\rightarrow$ LLaVA-OneVision)
has become the most important baseline series in the multimodal post-training field.

\evolution{Manually annotated multimodal instruction data}
  {GPT-4 automated generation + Two-Stage SFT}
  {LLaVA pipeline}

% --- InstructBLIP ---
\subsubsection{InstructBLIP: Instruction-Aware Visual Feature Extraction}
\label{sec:t5-instructblip}

\landmark{dai2023instructblip}
InstructBLIP \cite{dai2023instructblip} builds upon BLIP-2 \cite{li2023blip2},
and its core innovation is the \textbf{instruction-aware Q-Former} ---
introducing instruction information at the visual feature extraction stage,
enabling the Q-Former to extract \textit{task-relevant} visual features according to different task instructions.

\paragraph{Architecture.}
BLIP-2's Q-Former uses a fixed set of learned queries
to extract visual information from the frozen vision encoder's output via cross-attention.
InstructBLIP's key improvement is:
feeding instruction text tokens as additional input alongside the learned queries into the Q-Former,
so that the Q-Former's self-attention can simultaneously attend to both queries and instruction tokens,
while the cross-attention extracts information from visual features based on these instruction-conditioned queries.
Formally, let $\mathbf{Q}$ denote the learned queries, $\mathbf{I}$ the instruction tokens,
and $\mathbf{V}$ the visual features. Then:
\begin{equation}
  \mathbf{Q}' = \text{Q-Former}(\mathbf{Q}, \mathbf{I}, \mathbf{V})
\end{equation}
where self-attention operates over $[\mathbf{Q}; \mathbf{I}]$,
and cross-attention lets $\mathbf{Q}'$ attend to $\mathbf{V}$.
This design ensures that for the same image, different instructions guide the Q-Former to focus on different visual regions and attributes.

\paragraph{Training.}
InstructBLIP conducted the most systematic vision-language instruction tuning study to date:
constructing the training set from 26 public datasets covering 11 task categories
(image captioning, VQA, visual reasoning, image classification, etc.).
Each dataset is converted into instruction-following format using specially designed instruction templates.
During training, both the vision encoder and LLM remain frozen, and only the Q-Former is trained.

An important engineering detail is \textbf{balanced dataset sampling}:
directly sampling in proportion to dataset size would cause smaller datasets to be overwhelmed by larger ones.
InstructBLIP's strategy sets the sampling probability for each dataset to be proportional to the square root of its size:
$p_i \propto \sqrt{|D_i|}$, effectively balancing coverage across datasets.

\paragraph{Results.}
InstructBLIP achieved zero-shot SOTA across all 13 held-out (not used in training) datasets,
demonstrating the generalization capability of instruction-aware feature extraction.
Even more remarkably, the 4B-parameter InstructBLIP (Vicuna-7B backbone)
surpassed the 80B-parameter Flamingo on multiple benchmarks,
fully illustrating the efficiency advantage of instruction tuning:
\textit{high-quality post-training can compensate for an order-of-magnitude gap in parameter count}.

\begin{solutionbox}
\textbf{LLaVA vs InstructBLIP --- Two Visual Feature Extraction Paradigms:}
\begin{itemize}[leftmargin=2em]
  \item LLaVA: simple projection, feeding \textit{all} visual features directly into the LLM,
    letting the LLM's attention mechanism decide which visual information to focus on.
    Advantages: simple implementation, LLM receives complete visual information; Disadvantages: more visual tokens, higher computational cost.
  \item InstructBLIP: instruction-aware Q-Former,
    performing task-relevant filtering at the visual feature extraction stage.
    Advantages: fixed and fewer visual tokens (32 queries), strong task adaptability;
    Disadvantages: the Q-Former introduces additional training complexity and may discard visual information that is irrelevant to the current instruction but still valuable.
\end{itemize}
Subsequent developments show that LLaVA's simple projection approach (augmented with higher resolution and better projection designs)
has gradually become the mainstream choice, owing to its \textbf{simplicity} and \textbf{scalability}.
\end{solutionbox}

% ============================================================
% 5.3 Limitations & Branching
% ============================================================
\subsection{Limitations and Branching: Evolutionary Paths of Multimodal Alignment}
\label{sec:t5-branches}

Visual instruction tuning laid the foundational framework for multimodal post-training,
but in practical deployment, it exposed multiple pressing issues that gave rise to four branching evolutionary paths.

% --- Branch A: VLM Preference Optimization ---
\subsubsection{Branch A: VLM Preference Optimization}
\label{sec:t5-branch-a}

\begin{problembox}
SFT-trained VLMs exhibit alignment deficiencies similar to those of text-only LLMs:
models tend to generate responses that appear plausible but are unfaithful to image content,
and lack the ability to refuse uncertain questions.
How can preference optimization techniques be transferred from the text-only setting to the multimodal setting?
\end{problembox}

\paragraph{LLaVA-RLHF: The First Multimodal RLHF.}
\landmark{sun2023llava-rlhf}
LLaVA-RLHF \cite{sun2023llava-rlhf} is the first work to apply the complete RLHF pipeline to a VLM,
directly transferring the RLHF methodology from Thread~2
\crossref{2}{sec:t2}
to the multimodal setting.
Its training consists of three stages:

\begin{enumerate}[leftmargin=2em]
  \item \textbf{Multimodal SFT}:
    Building upon LLaVA, enhanced high-quality data
    (VQA-v2 \cite{goyal2017vqav2}, A-OKVQA \cite{schwenk2022aokvqa}, Flickr30k \cite{plummer2015flickr30k})
    is used for SFT, providing a better warm start.

  \item \textbf{Hallucination-Aware Human Preference Collection}:
    Approximately 10K human preference annotations are collected (costing approximately \$3,000).
    The annotation process is specifically designed to focus on hallucination:
    annotators are instructed to prioritize penalizing responses containing factual errors
    when comparing two responses,
    rather than judging solely based on fluency or level of detail.

  \item \textbf{Factually Augmented RLHF (Fact-RLHF)}:
    This is the most critical innovation of this work.
    In standard RLHF, the reward model receives only (image, question, response) as input,
    which is prone to \textbf{reward hacking} ---
    the reward model may assign high scores to ``confident but incorrect'' responses.
    Fact-RLHF's solution is to additionally append \textit{factual information}
    (detailed image captions and/or ground truth answers) to the reward model's input,
    enabling the reward model to verify the factual correctness of responses based on ground truth information.
    Formally, the reward model's input is extended from $(x_v, x_q, y)$ to $(x_v, x_q, y, x_{\text{fact}})$,
    where $x_{\text{fact}}$ is the factual augmentation.
\end{enumerate}

\noindent
On MMHal-Bench (a hallucination evaluation benchmark containing 96 image-question pairs
across 8 categories $\times$ 12 object topics),
LLaVA-RLHF achieved a 60\% reduction in hallucination compared to the LLaVA baseline.
On LLaVA-Bench, it attained 94\% of the GPT-4 evaluation score.

\paragraph{RLHF-V: Fine-Grained Correctional Feedback.}
RLHF-V \cite{yu2023rlhfv} proposes a more refined preference collection strategy:
annotators not only label preferred/dispreferred response pairs,
but also annotate, segment by segment, which specific portions of the dispreferred response contain hallucination,
and provide corresponding corrective text.
This \textit{fine-grained correctional feedback} enables preference learning
to focus on the genuinely problematic parts of a response, rather than making coarse-grained judgments about the entire response.

\paragraph{RLAIF-V: Open-Source AI Feedback.}
RLAIF-V \cite{yu2024rlaifv} further addresses the cost problem of human annotation:
it uses open-source VLMs (rather than GPT-4V) to automatically generate preference data,
forming a fully open-source AI feedback pipeline.
This work demonstrates the feasibility of AI feedback in VLM alignment,
echoing the idea of RLAIF \cite{lee2023rlaif} from Thread~2
\crossref{2}{sec:t2}.

\paragraph{mDPO: Conditional Preference Optimization.}
mDPO \cite{wang2024mdpo} modifies DPO for the multimodal setting.
The implicit assumption of standard DPO is that the difference in the reference policy's probabilities
for preferred and dispreferred responses primarily arises from textual preferences,
but in the multimodal setting, image conditioning significantly affects this distribution.
mDPO introduces an image-conditioned correction term,
enabling DPO's preference learning to better accommodate the characteristics of multimodal inputs.

\paragraph{LLaVA-Critic: Model as Evaluator.}
LLaVA-Critic \cite{xiong2024llava-critic} proposes training the VLM itself as an evaluator of multimodal outputs,
capable of scoring and ranking responses from itself or other models.
This work is directly aligned with the reward modeling ideas from Thread~2
\crossref{2}{sec:t2},
while also providing infrastructure for self-improvement and iterative DPO in the multimodal setting.

\paragraph{Other Works.}
Silkie \cite{li2023silkie} uses GPT-4V to perform preference ranking of responses from multiple VLMs,
constructing training data through preference distillation.
HA-DPO \cite{zhao2023hadpo} specifically constructs preference pairs targeting hallucination,
making DPO's optimization objective more focused on reducing hallucination.
SeVa \cite{zhu2024seva} proposes self-supervised visual preference alignment,
generating preference data by contrasting correct/incorrect visual inputs without requiring human annotation.
DRESS \cite{chen2023dress} introduces natural language feedback as an alignment signal,
going beyond simple preferred/dispreferred binary annotation.

\evolution{SFT-only VLM training}
  {Complete alignment pipeline with SFT $\rightarrow$ RLHF/DPO}
  {Multimodal preference optimization}

% --- Branch B: Hallucination Reduction ---
\subsubsection{Branch B: Hallucination Reduction}
\label{sec:t5-branch-b}

\begin{problembox}
VLMs frequently generate descriptions inconsistent with image content --- claiming the presence of objects that do not exist in the image (object hallucination),
or incorrectly describing object attributes and relationships.
How can multimodal hallucination be systematically evaluated and reduced?
\end{problembox}

\paragraph{Evaluation Benchmarks.}
\landmark{rohrbach2018chair}
CHAIR \cite{rohrbach2018chair} was the first to quantitatively evaluate object hallucination
by checking whether objects mentioned in captions exist in ground truth annotations,
but it covers only 80 COCO categories and is sensitive to prompts.
\landmark{li2023pope}
POPE \cite{li2023pope} transforms evaluation into a stable Yes/No binary classification,
using Random/Popular/Adversarial negative sampling strategies to reveal
the model's systematic hallucination tendency toward high-frequency and co-occurring objects,
as well as a strong bias toward ``Yes'' answers (Yes-ratio far exceeding 50\%).
MMBench \cite{liu2023mmbench} extends evaluation to 20+ fine-grained capability dimensions,
and HallusionBench \cite{guan2023hallusionbench} covers attribute, relation, and visual illusion scenarios.

\paragraph{Inference-Time Mitigation.}
VCD \cite{leng2023vcd} (CVPR 2024 Highlight) proposes training-free \textbf{Visual Contrastive Decoding}:
applying Gaussian noise perturbation to the input image to obtain a degraded input $v'$,
and subtracting its output (which reflects pure language priors and statistical biases) from the original output,
thereby amplifying the token probabilities that genuinely depend on visual signals.
Woodpecker \cite{yin2023woodpecker} constructs a training-free five-stage factual verification pipeline,
using Grounding DINO and BLIP-2 to perform claim-level verification of VLM outputs,
improving the accuracy of weaker models (mPLUG-Owl) by up to 24\%.

\paragraph{Training-Time Mitigation.}
LURE \cite{zhou2023lure} (ICLR 2024) identifies three factors causing hallucination ---
co-occurrence bias, decoding uncertainty, and position bias ---
and designs a post-processing revisor that reduces CHAIR$_S$ by 73.6\%.
HACL \cite{jiang2023hacl} uses contrastive learning
to push apart hallucinated and faithful representations in feature space.
HalluciDoctor \cite{yu2023hallucidoctor} eliminates hallucination from a data quality perspective, removing hallucination in the training data itself.
These training-time methods complement the DPO methods in Branch A ---
the former corrects data or representations, while the latter guides generation behavior through preference optimization.

\evolution{VLM evaluation without hallucination metrics}
  {Systematic evaluation (CHAIR/POPE) + dual-stage training/inference mitigation}
  {VCD/Woodpecker + LURE/HACL + DPO}

% --- Branch C: Grounding ---
\subsubsection{Branch C: Grounding and Fine-Grained Understanding}
\label{sec:t5-branch-c}

\begin{problembox}
Standard VLMs can only perform image-level understanding (``there is a cat in the image''),
but cannot perform region-level localization (``the cat is in the upper-left corner of the image [x1, y1, x2, y2]'').
How can VLMs be endowed with both grounding (text $\rightarrow$ location)
and referring (location $\rightarrow$ text) capabilities?
\end{problembox}

\paragraph{Kosmos-2: Pioneering Location Tokens.}
\landmark{peng2023kosmos2}
Kosmos-2 \cite{peng2023kosmos2} integrates spatial localization into the token vocabulary:
it introduces 1,024 location tokens and markup tokens,
embedding bounding boxes in hyperlink format within text.
This design unifies grounding and referring as standard next-token prediction,
without requiring additional detection heads or regression losses.
The accompanying GRIT dataset (91M images, 137M bounding boxes) provides the training foundation.

\paragraph{Shikra: Plain-Text Coordinate Scheme.}
Shikra \cite{chen2023shikra} (EMNLP 2024) demonstrates that spatial localization requires no architectural modifications ---
simply encoding coordinates as normalized decimal numbers (to 3 decimal places),
processed by the LLM's standard tokenizer, is sufficient to learn spatial reasoning through next-token prediction.
Shikra defines the \textbf{Referential Dialogue} (RD) format,
in which both users and models can embed coordinate references within natural language.
It achieves 87.83\% on RefCOCO val, demonstrating the competitiveness of the plain-text coordinate scheme.

\paragraph{Ferret: Arbitrary-Shape Regions.}
Ferret \cite{you2023ferret} (Apple, ICLR 2024) breaks beyond rectangular bounding box limitations,
using a \textbf{Spatial-Aware Visual Sampler} inspired by PointNet++
to support arbitrary-shape region inputs such as points, boxes, scribbles, and segmentation masks.
Combined with hybrid region representation and spatial hard negatives in training data,
Ferret-13B achieves 89.48\% on RefCOCO val.
LLaVA-Grounding \cite{zhang2023llava-grounding} jointly trains grounding with instruction tuning,
and Osprey \cite{yuan2023osprey} focuses on pixel-level understanding.

\evolution{Image-level understanding (whole-image QA)}
  {Region-level understanding (grounding + referring)}
  {Location tokens + Grounded data}

% --- Branch D: Multi-Image and Video ---
\subsubsection{Branch D: Multi-Image and Video Post-Training}
\label{sec:t5-branch-d}

\begin{problembox}
Early VLM post-training (LLaVA, InstructBLIP) focused on single-image understanding,
but real-world scenarios require understanding comparative relationships across multiple images or temporal dynamics in videos.
How can visual instruction tuning be extended from single images to multi-image and video settings?
\end{problembox}

\paragraph{Early Video-Language Models.}
\landmark{alayrac2022flamingo}
Flamingo's native support for interleaved image-text sequences provided an architectural reference for multi-image understanding.
VideoChat \cite{li2023videochat} explored two approaches: textualization (transcription via perception models) and end-to-end learning.
Video-LLaMA \cite{zhang2023videollama} (EMNLP 2023) constructs a dual-branch architecture for visual and audio,
using a Video Q-Former to aggregate temporal features and an Audio Q-Former to process audio via ImageBind cross-modal alignment.
Video-ChatGPT \cite{maaz2023videochatgpt} proposes temporal-spatial dual-pooling
and constructs a 100K video instruction tuning dataset.
These early works validated the feasibility of video instruction tuning,
but were limited in long video and deep temporal reasoning.

\paragraph{LLaVA-OneVision: Unified Architecture Convergence.}
LLaVA-OneVision \cite{li2024llava-onevision} is the first open-source model to
simultaneously achieve SOTA across single-image, multi-image, and video tasks.
Its \textbf{AnyRes} strategy segments high-resolution images into sub-image tiles that are independently encoded and then concatenated with a global overview.
Three-stage training (alignment $\rightarrow$ knowledge learning $\rightarrow$ instruction tuning)
enables video understanding capabilities to naturally emerge through joint training (image-to-video transfer), without requiring dedicated video pre-training.
The 72B version achieves 91.3\% on DocVQA (vs.\ GPT-4V 88.4\%),
77.6\% on multi-image Mantis (vs.\ GPT-4V 62.7\%), and 66.2\% on video VideoMME (vs.\ GPT-4V 59.9\%).
LLaVA-NeXT-Interleave \cite{li2024llava-next-interleave}
further extends to interleaved image-text inputs.

\paragraph{LLaVA-Video: Synthetic Video Instruction Data.}
LLaVA-Video \cite{zhang2024llava-video} constructs the LLaVA-Video-178K dataset (1.3M instruction samples)
through GPT-4o hierarchical caption generation (1 FPS high sampling rate, three-level progressive summarization).
It introduces \textbf{SlowFast} dual-path video encoding (Slow preserves spatial details, Fast processes more frames),
with the 72B model achieving 70.5\%/76.9\% on VideoMME (without/with subtitles).
LLaVA-Hound-DPO \cite{zhang2024llava-hound} is the first to apply DPO to video VLMs.

\evolution{Single-image Visual Instruction Tuning}
  {Unified training across single-image + multi-image + video}
  {LLaVA-OneVision pipeline}

% ============================================================
\subsection{Convergence and Frontiers: Formation of the Standard Pipeline}
\label{sec:t5-convergence}

After the exploratory branching evolution phase, the multimodal post-training field exhibits clear convergence trends along several dimensions.

\paragraph{Convergence Trend 1: SFT + DPO as the Standard Pipeline.}
Similar to how text-only alignment converged from RLHF to DPO in Thread~2
\crossref{2}{sec:t2},
multimodal post-training has also progressively converged from complex PPO-based RLHF (such as LLaVA-RLHF)
to the simpler SFT + DPO pipeline.
The three-stage training of LLaVA-OneVision \cite{li2024llava-onevision}
(Stage 1: alignment pre-training; Stage 2: multi-task SFT; Stage 3: DPO)
represents the latest manifestation of this convergence.
DPO's success in multimodal settings is attributable to its implementation simplicity
and direct effectiveness for hallucination reduction.

\paragraph{Convergence Trend 2: Diversity and Fusion of Vision Encoders.}
Cambrian-1 \cite{tong2024cambrian} systematically investigated the impact of visual representation on VLM performance.
Its core finding is that
no single visual encoder is optimal across all tasks ---
CLIP excels at semantic understanding, DINOv2 excels at spatial/geometric features,
and SigLIP outperforms CLIP in certain scenarios.
Cambrian-1 therefore proposes the \textbf{Spatial Vision Aggregator (SVA)},
which aggregates features from multiple heterogeneous visual encoders through learnable spatial tokens,
achieving multi-encoder fusion.
This work reveals an important insight:
\textit{the capability bottleneck of VLMs may not lie in the LLM or the post-training method,
but in the visual representation itself}.

\paragraph{Convergence Trend 3: Unification of Single-Image, Multi-Image, and Video.}
Early works tended to build independent models and training pipelines for different modalities
(LLaVA for single images, VideoChat for video),
but the latest trend is to cover all modalities within a unified architecture and training framework.
LLaVA-OneVision \cite{li2024llava-onevision} and
Qwen-VL \cite{bai2023qwen-vl} both demonstrate the viability of this unified approach.

\paragraph{Current Frontier.}
Synthesizing the above convergence trends, the current frontier of multimodal post-training can be summarized as:
\textit{Multi-encoder visual backbone $+$ high-resolution processing $+$
multi-task SFT on diverse data $+$ DPO for hallucination reduction $+$
unified single/multi-image/video support}.
This paradigm is embodied in works such as Cambrian-1 and LLaVA-OneVision.

\subsubsection{2025: Multimodal Reasoning via RL}
\label{sec:t5-reasoning-rl}

In 2025, the success of the RLVR paradigm in Thread~4 (\crossref{4}{sec:t4-thinking-models})
rapidly spilled over into the multimodal domain,
marking a paradigm shift from ``SFT + DPO'' to ``SFT + RL for reasoning.''

\paragraph{Rule-Based RL for VLM.}
R1-VL \cite{zhang2025r1vl} (ICCV 2025) extends step-wise GRPO to visual mathematical reasoning,
validating the effectiveness of RLVR in multimodal settings.
MM-Eureka \cite{meng2025mmeureka} and VLM-R1 \cite{shen2025vlmr1}
advance R1-style VLMs in mathematical/scientific reasoning and training stability, respectively.
MM-RLHF \cite{zhang2025mmrlhf} (ICML 2025) systematically extends RLHF to multimodal settings,
constructing large-scale multimodal preference datasets and multi-dimensional reward models.
TLDR \cite{fu2025tldr} refines rewards from sequence-level to token-level,
and Multimodal RewardBench \cite{yasunaga2025mmrewardbench} provides a systematic evaluation benchmark.

\paragraph{URSA: Process-Supervised RL for Multimodal Math.}
Luo et al.\ \cite{luo2025ursa}'s \textbf{URSA} (NeurIPS 2025) constructs a three-stage pipeline:
MMathCoT-1M (large-scale multimodal chain-of-thought data),
DualMath-1.1M (process-level supervision covering both logical correctness and perceptual consistency),
and \textbf{PS-GRPO} (Process-Supervised GRPO) ---
the first method to apply PRM-guided online RL to multimodal reasoning.
URSA-8B surpasses Gemma3-12B and GPT-4o on average across 6 mathematical benchmarks
(by 8.4\% and 2.7\%, respectively),
directly extending the process reward ideas from Thread~4
\crossref{4}{sec:t4-thinking-models}
to vision-language joint reasoning settings.

\paragraph{VisualPRM: Multimodal Process Reward Model.}
The InternVL team's \textbf{VisualPRM} \cite{wang2025visualprm} (8B parameters)
scores \textit{each step} of vision-language chain-of-thought rather than evaluating only the final answer.
Through an automated pipeline, it constructs the VisualPRM400K training dataset and releases the VisualProcessBench evaluation benchmark.
Under best-of-N sampling, it improves InternVL2.5-78B by an average of 5.9 points and the 8B model by 8.4 points,
providing open-source infrastructure for VLM test-time scaling.

\paragraph{MVoT: Thinking in Images.}
\textbf{MVoT} \cite{li2025mvot} (ICML 2025) proposes a qualitatively different reasoning paradigm:
interleaving the generation of image tokens and text tokens within the chain-of-thought,
allowing the model to ``think in images'' rather than ``think about images.''
It surpasses text-only CoT in the most challenging dynamic spatial reasoning scenarios,
demonstrating that certain reasoning tasks genuinely require visual intermediate representations.

\paragraph{Paradigm Shift Summary.}
o3/o4-mini \cite{openai2025o3} possesses native multimodal reasoning capabilities, further reinforcing this trend.
In 2025, multimodal reasoning advances along three directions:
\textit{refinement of training signals} (process-level reward),
\textit{test-time compute scaling} (PRM-guided search),
and \textit{expansion of reasoning modalities} (from text-only CoT to visual CoT).
The frontier of multimodal post-training is evolving from
\textit{SFT + DPO for hallucination reduction}
to \textit{SFT + RL for reasoning-capable VLMs}.

\subsubsection{2025: VLM Architecture Engineering and Open-Source Scaling}
\label{sec:t5-scaling}

In 2024--2025, open-source VLMs such as InternVL, Qwen-VL, and DeepSeek-VL have approached or even surpassed GPT-4o,
systematically answering key engineering questions in VLM post-training.

\paragraph{Dynamic Resolution: From Fixed to Native Resolution.}
The shift from fixed resolution to native resolution is the consensus of 2024--2025:
LLaVA-OneVision \cite{li2024llava-onevision}'s \textbf{AnyRes}
segments images into tiles that are independently encoded and then compressed via spatial pooling,
becoming the most widely adopted strategy.
Qwen2-VL \cite{wang2024qwen2vl}'s Naive Dynamic Resolution
directly processes original resolutions and introduces \textbf{M-RoPE}
to unify temporal-spatial positional encoding.
Qwen2.5-VL \cite{bai2025qwen25vl} introduces Window Attention within the ViT
(with only 4 layers of full attention), significantly reducing computational overhead,
while treating precise object localization and document parsing as first-class post-training objectives.

\paragraph{Vision Encoder Scaling.}
Cambrian-1 \cite{tong2024cambrian} systematically revealed the critical importance of visual representation
(see \S\ref{sec:t5-convergence}).
\textbf{InternVL2.5} \cite{chen2024internvl25}'s 6B InternViT significantly reduced the dependence on data volume ---
achieving superior results with only one-tenth of Qwen2-VL-72B's training tokens,
revealing an underappreciated empirical law:
\textit{investing in the vision encoder can substitute for investing in data}.
\textbf{InternVL3} \cite{zhu2025internvl3} proposes Native Multimodal Pre-Training ---
interleaving multimodal data with text during the pre-training stage,
challenging the fundamental assumption that ``multimodal capabilities come from the adaptation stage.''

\paragraph{Efficient Architecture and Data Pipelines.}
\textbf{DeepSeek-VL2} \cite{wu2024deepseekvl2} introduces MoE + MLA into VLMs,
surpassing GPT-4o on OCRBench with fewer activated parameters (834 vs.\ 736).
\textbf{Molmo} \cite{deitke2024molmo} (CVPR 2025) offers a unique perspective from the data side:
the PixMo dataset is entirely constructed through human annotation (verbal descriptions + 2D pointing), without using VLM-synthesized data.
Molmo-72B is second only to GPT-4o, demonstrating that
\textit{eliminating the circular dependency on synthetic data} has a first-order impact on post-training quality.

\begin{solutionbox}
\textbf{Post-Training Insights from the Open-Source VLM Wave:}
\begin{itemize}[leftmargin=2em]
  \item \textbf{Resolution}: Dynamic tiling + spatial pooling has become the de facto standard.
  \item \textbf{Vision Encoder}: Large-scale ViTs and ViTs trained from scratch
    both outperform directly reusing CLIP/SigLIP; the vision encoder is the true bottleneck.
  \item \textbf{Data quality > scale}: Molmo's purely human-annotated data outperforms competitors using massive synthetic data,
    and InternVL2.5 surpasses Qwen2-VL with one-tenth the tokens.
\end{itemize}
\end{solutionbox}

\subsubsection{2025: Unified Multimodal Generation}
\label{sec:t5-unified-generation}

All the VLM works discussed above --- from LLaVA to InternVL3 ---
share a fundamental assumption: \textbf{the model only performs understanding, not generation}.
They receive images and output text, but cannot generate images or videos.
In 2024--2025, a new paradigm shift is underway:
\textbf{unifying visual understanding and visual generation within the same model and training framework}.
This is not only an architectural challenge but also poses entirely new requirements for post-training ---
how can understanding quality and generation quality be simultaneously optimized within the same alignment pipeline?

\paragraph{Architectural Spectrum: From Early Fusion to Decoupled Encoding.}
The architectural designs for unified multimodal generation form a clear spectrum:

\begin{itemize}[leftmargin=2em]
  \item \textbf{Early Fusion --- Chameleon} \cite{team2024chameleon}:
    A 7B/34B model proposed by Meta,
    which from pre-training onward treats images and text as the same discrete token vocabulary,
    trained on 4.4T mixed-modal tokens.
    Chameleon is the first large-scale early-fusion token-based foundation model,
    eliminating the ``late fusion'' mode of adding vision adapters onto a frozen LLM.
    The stability challenges encountered during training (resolved through QK-Norm and z-loss)
    provided important experience for subsequent unified model training.

  \item \textbf{Pure Autoregressive --- Emu3} \cite{wang2024emu3}:
    Tokenizes images, text, and video into a unified discrete space,
    training a single transformer using only next-token prediction.
    No diffusion, no separate encoder, no compound loss functions.
    Emu3 simultaneously surpasses SDXL (text-to-image), LLaVA-1.6 (understanding),
    and OpenSora-1.2 (video generation).
    Published in Nature \cite{wang2024emu3},
    this work represents the most extreme position on the unified architecture spectrum ---
    if autoregressive prediction is sufficient,
    then the entire suite of text LLM post-training techniques (SFT/RLHF) can be directly transferred.

  \item \textbf{Hybrid AR + Diffusion --- Show-o / Transfusion}:
    Show-o \cite{xie2024showo} (ICLR 2025) uses autoregressive modeling for text and discrete diffusion for images
    within the same transformer,
    handling mixed tasks such as VQA, text-to-image, and inpainting in a single forward pass.
    Transfusion \cite{zhou2024transfusion} (Meta) goes further,
    applying next-token prediction loss to text
    and diffusion loss to \textit{continuous} image patches,
    without needing to discretize images into tokens.
    These two works represent a pragmatic middle ground ---
    acknowledging that different modalities may require different training objectives, but still achieving this within a unified architecture.

  \item \textbf{Decoupled Encoding --- Janus / Janus-Pro} \cite{chen2024janus,chen2025januspro}:
    Janus (CVPR 2025), proposed by DeepSeek, identifies a fundamental tension:
    the visual representation needed for understanding (rich semantic features, CLIP-style)
    and the visual representation needed for generation (fine-grained spatial features for decoding)
    are \textit{structurally different}.
    Janus's solution is to use \textbf{two separate visual encoding paths} ---
    one for understanding and one for generation ---
    while sharing the same LLM backbone and autoregressive framework.
    Janus-Pro scales this design to 1B and 7B parameters,
    surpassing prior unified models in both understanding and generation.
\end{itemize}

\paragraph{New Challenges for Post-Training.}
Unified generation models introduce fundamentally new problems for post-training:
\begin{itemize}[leftmargin=2em]
  \item \textbf{Multi-Objective Alignment}:
    How can understanding accuracy (faithfulness, reduced hallucination)
    and generation quality (aesthetic quality, factual consistency, user intent matching)
    be simultaneously optimized within the same RLHF/DPO pipeline?
    The two objectives may conflict --- training signals beneficial for understanding may not be beneficial for generation.
  \item \textbf{Tokenization Determines Post-Training Feasibility}:
    Emu3's pure autoregressive approach means that the entire suite of text LLM post-training techniques
    can be directly transferred;
    whereas Transfusion's hybrid loss design requires designing different reward signals
    and preference learning strategies for different modalities.
  \item \textbf{Exponential Growth in Data Requirements}:
    Understanding and generation each require high-quality data,
    and joint training of both may require even more interleaved data
    to prevent catastrophic forgetting between capabilities.
\end{itemize}

\evolution{Understanding-only VLMs (receive images, output text)}
  {Unified Understanding + Generation (same model understands and generates images/video)}
  {Decoupled / AR / Hybrid architectural paths}

\subsubsection{2025: Omni-Modal Post-Training}
\label{sec:t5-omni}

Unified generation merges visual understanding and visual generation into the same model,
while \textbf{omni-modal} models extend this unification to \textit{all modalities} ---
end-to-end joint processing of text, vision, audio, and even speech.
The core post-training challenge of this direction lies in the fact that
different modalities have fundamentally different \textit{time scales} (audio in milliseconds, video in frames, text in tokens)
and \textit{interaction modes} (speech requires real-time duplex, text can be asynchronous).
How can these differences be coordinated within a single alignment pipeline?

\paragraph{GPT-4o: Commercial Validation of the Omni Paradigm.}
OpenAI's GPT-4o \cite{openai2024gpt4o} is the first end-to-end trained omni model,
accepting any combination of text, audio, image, and video as input,
and generating text, audio, and (since March 2025) image output.
Audio response latency is approximately 320ms, approaching human conversational speed.
GPT-4o eliminates the traditional cascade pipeline (ASR $\rightarrow$ LLM $\rightarrow$ TTS),
demonstrating the feasibility of end-to-end omni training at industrial scale,
and immediately triggering a competition in the open-source community.

\paragraph{Qwen2.5-Omni: Thinker-Talker Architecture.}
The Qwen team proposes \textbf{Qwen2.5-Omni} \cite{xu2025qwen25omni},
introducing the \textbf{Thinker-Talker} dual-module architecture:
the Thinker (LLM) is responsible for multimodal reasoning, and the Talker handles streaming speech generation.
The key technical innovation is \textbf{TMRoPE} (Time-aligned Multimodal Rotary Position Embedding),
which synchronizes time-sensitive audio and video inputs at the positional encoding level,
solving the alignment problem for multi-stream temporal data.
A sliding-window DiT decoder in the Talker module handles streaming audio token generation.
On speech instruction following benchmarks (MMLU, GSM8K),
Qwen2.5-Omni's speech input performance \textit{matches for the first time} its text input performance ---
a significant milestone for omni models.

\paragraph{VITA: Open-Source Omni with Duplex Interaction.}
\textbf{VITA} \cite{fu2024vita} is the first open-source omni multimodal LLM,
supporting simultaneous processing of video, image, text, and audio,
and featuring \textbf{duplex interaction} capability --- the model can be interrupted by the user while generating a response.
The subsequent VITA-1.5 (NeurIPS 2025) reduced end-to-end speech interaction latency
from approximately 4 seconds to approximately 1.5 seconds, with multimodal benchmark average scores improving from 59.8 to 70.8.
Duplex interaction introduces post-training challenges that standard instruction-following pipelines
have never addressed:
How to generate output while continuously receiving input?
How to train the model to appropriately interrupt its own responses?

\paragraph{Phi-4-Multimodal: Mixture-of-LoRAs.}
Microsoft's \textbf{Phi-4-Multimodal} \cite{abouelenin2025phi4mini} (5.6B parameters)
proposes a parameter-efficient omni adaptation scheme:
\textbf{Mixture-of-LoRAs} --- training independent LoRA adapters for each modality (speech, vision, text),
dynamically activated based on input type.
This design minimizes cross-modal interference while sharing the LLM backbone.
Phi-4-Multimodal is trained on 5T text tokens, 2.3M hours of speech, and 1.1T image-text tokens,
surpassing larger specialized vision-language and speech-language models with only 5.6B parameters.
Mixture-of-LoRAs is an efficient alternative to full fine-tuning,
directly connected to the parameter-efficient post-training methodology in Thread~7
\crossref{7}{sec:t7}.

\begin{solutionbox}
\textbf{Three Paradigms for Omni-Modal Architecture Design:}
\begin{itemize}[leftmargin=2em]
  \item \textbf{End-to-end monolithic} (GPT-4o):
    All modalities are trained end-to-end within the same model; highest quality but greatest training cost.
  \item \textbf{Modular decomposition} (Qwen2.5-Omni's Thinker-Talker):
    Reasoning and real-time generation are decoupled into independent modules,
    achieving a balance between quality and latency.
  \item \textbf{Adapter-based} (Phi-4-Multimodal's MoE-LoRA):
    The LLM backbone is kept unchanged, and capabilities are injected through modality-specific adapters;
    most efficient but may limit the depth of cross-modal interaction.
\end{itemize}
\end{solutionbox}

\evolution{Text + Vision bimodal VLM}
  {Text + Vision + Audio omni-modal Omni Model}
  {Thinker-Talker, Mixture-of-LoRAs, End-to-End Training}

\subsubsection{2025--2026: Emerging Frontiers}
\label{sec:t5-emerging}
\label{sec:t5-self-evolving}

The aforementioned RLVR and test-time scaling methods have opened multiple frontier directions.
The latest works in 2025--2026 advance the boundaries of VLM post-training along six paths.

\paragraph{Self-Evolving VLM Training.}
Static datasets face the problem of training saturation --- when the model's pass rate approaches 0 or 1, the gradient signal vanishes.
Jia et al.\ \cite{jia2026dpe}'s \textbf{DPE} (Diagnostic-driven Progressive Evolution)
reformulates training as a Diagnose $\rightarrow$ Produce $\rightarrow$ Enhance closed loop:
diagnosing the model's capability profile, automatically generating training data targeting weak dimensions, and reinforcing through GRPO.
The key $p(1-p)$ analysis shows that difficulty-aware filtering has an information-theoretic foundation ---
gradient signal variance is maximized at the model's decision boundary ($p=0.5$).
Merely 3K DPE samples surpass the effectiveness of 47K static data samples.
ProbeLLM \cite{huang2026probellm} uses MCTS to systematically probe model failure modes,
complementing DPE's diagnostic step
(see Chapter~\ref{sec:t8}, Directions 7--9).

\paragraph{Test-Time Compute Scaling.}
\label{sec:t5-test-time-compute}
VLM test-time scaling is far less stable than LLM test-time scaling:
visual perception errors are global --- once the model ``misreads'' an image, subsequent reasoning cannot self-correct.
VisVM \cite{wang2024visvm} trains a dedicated vision value model to guide beam search away from hallucination paths.
VL-PRM \cite{ong2025vlprm} transfers the process reward model to multimodal settings,
introducing perception-level supervision that simultaneously checks the correctness of both logic and visual grounding,
complementing URSA's PS-GRPO in \S\ref{sec:t5-reasoning-rl}.
RD-VLA \cite{tur2026rdvla} proposes iterative reasoning in latent space rather than token-space search,
achieving over 90\% success rate after 4 iterations on tasks with 0\% success in a single iteration, at 80$\times$ faster speed.
Empirical analysis \cite{ahmadpour2025ttsvisonlanguage} confirms that VLM test-time scaling
is nearly ineffective on perception-dominated tasks, with benefits highly dependent on task type.

\paragraph{Multimodal Reward Models.}
\label{sec:t5-reward-models}
Multimodal reward modeling has exposed a disturbing reality:
even the strongest VLMs serving as judges achieve only approximately 65--72\% accuracy
\cite{li2024vlrewardbench,yasunaga2025mmrewardbench}.
The bottleneck is not reasoning capability but perceptual reliability ---
a judge that ``cannot see the image clearly'' cannot reliably assess the faithfulness of responses.
MM-RLHF \cite{zhang2025mmrlhf} provides 120K fine-grained multimodal preference data;
Skywork-VL Reward \cite{wang2025skyworkvlreward} validates that
dedicated reward model training is more effective than using general VLMs as judges ---
consistent with the evolution of text-only reward models in Thread~2
(\crossref{2}{sec:t2}).
Critic-V \cite{zhang2024criticv} and Judge-MCTS \cite{chen2026mjudgebench}
advance the evolution from ``black-box scoring'' to ``structured critique.''

\paragraph{Domain-Specialized VLM RL.}
\label{sec:t5-domain-rl}
GRPO is expanding into three high-value vertical domains:
SCoPE VLM \cite{lim2025scopevlm} uses Episodic GRPO to train scroll-or-answer decisions in document navigation;
BigCharts-R1 \cite{masry2025bigchartsr1} designs chart-specific rewards covering numerical estimation, computation, and visual reasoning;
UI-TARS-2 \cite{wang2025uitars2} advances VLM RL from single-turn QA to multi-turn GUI interaction ---
through a Data Flywheel (alternating continual SFT + RL iterations)
and a stabilized multi-turn RL framework,
achieving 47.5\% on OSWorld, significantly surpassing Claude and OpenAI agent baselines.

\paragraph{Spatial \& 3D Reasoning.}
\label{sec:t5-spatial-reasoning}
Quantitative spatial reasoning (distance, size, direction) is a systematic blind spot of VLMs.
SpatialVLM \cite{chen2024spatialvlm} (CVPR 2024) automatically generates 2 billion spatial VQA samples from 10 million images,
revealing that ``capability deficiency is often data deficiency, not architectural deficiency.''
MetaSpatial \cite{pan2025metaspatial} is the first to introduce RL to 3D spatial reasoning
(3D Spatial Policy Optimization).
Smooth Operator \cite{jiao2026smoothoperator} discovers that
VLMs may already possess latent 3D reasoning capabilities,
requiring only RL post-training on 50K verifiable data samples to ``activate'' them ---
echoing DeepSeek-R1's finding in Thread~4 that ``RL awakens dormant capabilities''
(\crossref{4}{sec:t4-thinking-models}).

\paragraph{World Models as Data Engines.}
\label{sec:t5-world-models}
As video generation models mature,
a new paradigm is emerging: using world models to generate synthetic data for VLM/VLA post-training.
NVIDIA's Cosmos platform \cite{agarwal2025cosmoswfm,ali2025cosmospredict25} builds world model infrastructure,
and Cosmos-Reason1 \cite{azzolini2025cosmosreason1} is the first to systematically treat physical common sense as a post-training objective.
GigaBrain/GigaWorld \cite{ye2025gigabrain0,ye2025gigaworld0}
demonstrate that VLAs post-trained entirely on synthetic data can achieve strong generalization on real robot tasks.
VLAW \cite{guo2026vlaw} proposes iterative co-improvement between VLA policies and world models,
upgrading the world model from a data generator to a virtual training environment.
DPE generates questions from real images, while world models generate images from physical simulations ---
the fusion of both is a direction that remains to be explored
(see the future directions discussion in Chapter~\ref{sec:t8}).

\begin{evolutionbox}
\textbf{Six Frontier Directions in VLM Post-Training}

Static RLVR $\;\xrightarrow{\text{diagnostic-driven}}\;$ Self-evolving training (DPE)
\qquad
General GRPO $\;\xrightarrow{\text{task-specific reward}}\;$ Domain-specialized RL (document/chart/GUI)
\qquad
2D understanding $\;\xrightarrow{\text{spatial data + RL}}\;$ 3D spatial reasoning post-training
\qquad
Real data $\;\xrightarrow{\text{world model}}\;$ Synthetic data-driven VLA training
\end{evolutionbox}

% ============================================================
% 5.5 Open Problems
% ============================================================
\subsection{Open Problems}
\label{sec:t5-open}

Despite significant progress, multimodal post-training still faces several unresolved core problems.
Note that some of the following problems have been partially mitigated by the works discussed above,
but fundamental challenges remain unsolved.

\begin{openproblembox}
\textbf{Open Problem 1: Fine-Grained Visual Understanding.}

Qwen2.5-VL's bounding box/point localization capabilities and Molmo's pointing dataset
have advanced fine-grained understanding to a practical level.
However, for tasks requiring \textit{pixel-level} precise understanding (such as medical image analysis, remote sensing image interpretation,
and industrial defect detection), the granularity of existing methods remains insufficient.
The core bottleneck is that
the visual token granularity of current VLMs (typically 14$\times$14 pixels per token)
exhibits an order-of-magnitude resolution gap relative to the requirements of pixel-level tasks.
How to bridge this gap during the post-training stage ---
whether by introducing segmentation-aware training data or designing multi-scale visual token mechanisms ---
remains an open problem.
\end{openproblembox}

\begin{openproblembox}
\textbf{Open Problem 2: Long Video Understanding \& Temporal Alignment.}

Multiple mitigation approaches have appeared in 2025:
Video-XL \cite{shu2024videoxl} (CVPR 2025) proposes \textbf{Visual Summarization Tokens (VST)} ---
periodically inserting special tokens into the video token sequence,
where the model learns to compress interval information into these tokens, achieving a 16$\times$ compression ratio,
enabling a single A100 GPU to process thousands of frames in long videos.
VideoLLaMA 3 \cite{zhang2025videollama3} proposes an ``images first, video second'' training philosophy ---
first building a strong visual foundation on high-quality image-text data, then extending to video,
with similarity-based frame merging to reduce redundant tokens.
MLVU \cite{zhou2024mlvu} (CVPR 2025) establishes the first standard evaluation benchmark
covering 9 task categories for videos ranging from 3 minutes to 2 hours,
revealing that even the strongest models (GPT-4o) achieve only 64.6\% on multiple-choice tasks.

Nevertheless, \textbf{hour-scale temporal grounding} (localizing precise time intervals of events)
and \textbf{streaming video understanding} (performing real-time reasoning while video input continues)
remain unsolved challenges.
InternLM-XComposer2.5-OmniLive \cite{chen2024omnilive}
separates perception, long-term memory, and reasoning into three independent modules,
representing an initial exploration in the streaming direction,
but systematic methods for effectively training ``memory management'' capabilities during post-training are still lacking.
\end{openproblembox}

\begin{openproblembox}
\textbf{Open Problem 3: Omni-Modal Safety.}

The visual modality opens new attack surfaces for jailbreaking attacks.
Qi et al.\ \cite{qi2023visual} demonstrated that embedding adversarial perturbations into images
can successfully bypass the text-based safety alignment of VLMs;
FigStep \cite{gong2023figstep} showed that simply rendering harmful instructions into images
can bypass text safety filters --- a simpler attack;
MM-SafetyBench \cite{liu2023mm-safetybench} systematically evaluated the effectiveness of various attack methods;
and VLGuard \cite{zong2024vlguard} proposed a method for safety fine-tuning at extremely low cost.

With the emergence of omni-modal models (\S\ref{sec:t5-omni}),
the attack surface for safety concerns expands further ---
the audio modality provides yet another unguarded backdoor
(e.g., bypassing safety filters through covert instructions embedded in speech).
Current VLM safety alignment is primarily built for the text modality
\crossref{3}{sec:t3},
and how to construct \textbf{modality-agnostic} safety alignment methods,
enabling models to possess robust refusal capabilities against harmful inputs from \textit{any} modality,
is a problem of increasing urgency as the number of modalities grows.
\end{openproblembox}

\begin{openproblembox}
\textbf{Open Problem 4: Deep Multimodal Reasoning.}

URSA's PS-GRPO and VisualPRM have introduced process-level rewards
into multimodal reasoning training (\S\ref{sec:t5-reasoning-rl}),
and MVoT has explored the new reasoning modality of ``thinking in images.''
However, current multimodal reasoning remains heavily dependent on \textit{text-based chain-of-thought} ---
even when addressing visual problems, the reasoning chain is purely textual,
with visual information processed only once at the input stage and then ``forgotten'' in the depths of attention.

Truly deep multimodal reasoning should be able to
\textit{repeatedly consult and reinterpret visual information} during the reasoning process
(similar to how humans repeatedly refer to diagrams when solving problems).
How to train this ``visual look-back'' capability during the post-training stage ---
whether through MVoT-style visual intermediate steps,
or through a paradigm similar to agentic tool use (\crossref{6}{sec:t6})
where the vision encoder is ``called as a tool'' ---
remains an open research direction.
\end{openproblembox}

\begin{openproblembox}
\textbf{Open Problem 5: Post-Training for Unified Generation.}

The unified models in \S\ref{sec:t5-unified-generation} (Janus, Emu3, Show-o, Transfusion)
have demonstrated the architectural feasibility of simultaneously achieving understanding and generation.
However, the \textbf{post-training methodology} severely lags behind architectural innovation ---
these models currently rely primarily on SFT,
and virtually no work has systematically studied how to perform preference optimization on unified models.
The core difficulty lies in \textbf{multi-objective alignment}:
the alignment objectives for understanding tasks (faithfulness, reduced hallucination)
and the alignment objectives for generation tasks (aesthetic quality, user intent matching)
may conflict.
How to design reward models and preference learning strategies
that can simultaneously optimize both types of objectives
is a critical bottleneck for unified models to progress from ``capability demonstration'' to ``production-grade deployment.''
\end{openproblembox}

\begin{openproblembox}
\textbf{Open Problem 6: Cross-Modal Consistency in Omni-Modal Alignment.}

Omni models (\S\ref{sec:t5-omni}) require the model to provide consistent responses
to the same query across different input modalities ---
asking ``What is the capital of Paris?'' via text and asking the same question via speech
should not produce answers of different quality.
Qwen2.5-Omni's speech-text performance matching is a milestone,
but deeper consistency issues have yet to be addressed:
when the model simultaneously receives conflicting multimodal signals
(e.g., the image shows ``red'' but the speech says ``blue''),
how can clear modality priority and conflict resolution mechanisms be established through post-training?
Furthermore, duplex interaction (VITA) requires the model to
continuously process new inputs while generating output and adjust its responses in a timely manner ---
a training scenario entirely unaddressed by standard SFT / DPO pipelines.
\end{openproblembox}

\subsection{Chapter Summary}

\begin{figure}[!ht]
\centering
\begin{tikzpicture}[
  node distance=0.6cm and 1.2cm,
  every node/.style={font=\small},
  milestone/.style={rectangle, rounded corners, draw=thread5, fill=thread5!10,
    text width=4.2cm, minimum height=0.8cm, align=center, font=\small\bfseries},
  branch/.style={rectangle, rounded corners, draw=gray!60, fill=gray!5,
    text width=3.2cm, minimum height=0.7cm, align=center, font=\small},
  arrow/.style={-{Stealth[length=2.5mm]}, thick, thread5!70},
  brancharrow/.style={-{Stealth[length=2mm]}, gray!60, thick},
]
% Main timeline
\node[milestone] (flamingo) {Flamingo (2022)\\Visual In-Context Learning};
\node[milestone, below=1.2cm of flamingo] (llava) {LLaVA (2023)\\Visual Instruction Tuning};
\node[milestone, below=1.2cm of llava] (rlhf) {LLaVA-RLHF (2023)\\Multimodal Preference Opt.};

% Branches
\node[branch, below left=1.2cm and 0.5cm of rlhf] (brA) {Branch A: Architecture\\InstructBLIP, Kosmos-2};
\node[branch, below right=1.2cm and 0.5cm of rlhf] (brB) {Branch B: Hallucination\\POPE, RLHF-V, mDPO};

% Convergence
\node[milestone, below=2.8cm of rlhf] (onevision) {LLaVA-OneVision (2024)\\Unified Multi-task};
\node[milestone, below=1.2cm of onevision] (mmreason) {Multimodal Reasoning\\(2025)};
\node[branch, below=1.2cm of mmreason] (conv) {Convergence: From ``seeing''\\to ``reasoning'' via CoT + Vision};

% Arrows
\draw[arrow] (flamingo) -- (llava);
\draw[arrow] (llava) -- (rlhf);
\draw[brancharrow] (rlhf) -- (brA);
\draw[brancharrow] (rlhf) -- (brB);
\draw[arrow] (rlhf) -- (onevision);
\draw[arrow] (onevision) -- (mmreason);
\draw[brancharrow] (brA) -- (onevision);
\draw[brancharrow] (brB) -- (onevision);
\draw[arrow] (mmreason) -- (conv);
\end{tikzpicture}
\caption{Thread 5 evolutionary path: from visual in-context learning to multimodal reasoning integration.
Orange nodes represent landmark works; gray nodes represent important branches.}
\label{fig:multimodal-evolution}
\end{figure}

\begin{table}[!ht]
\centering
\caption{Comparison of core methods in Thread 5. Vision-LLM Bridge refers to the mechanism for connecting visual information to the LLM.}
\label{tab:multimodal-comparison}
\small
\begin{tabularx}{\textwidth}{l >{\raggedright\arraybackslash}p{2.5cm} >{\raggedright\arraybackslash}p{2.5cm} >{\raggedright\arraybackslash}X}
\toprule
\textbf{Method} & \textbf{Vision-LLM Bridge} & \textbf{Post-training} & \textbf{Key Contribution} \\
\midrule
Flamingo & Gated cross-attention & Frozen LLM; visual SFT & Visual in-context learning (few-shot) \\
LLaVA & Linear projection & Two-stage SFT (align + instruct) & Visual instruction tuning at $<$\$300 \\
InstructBLIP & Instruction-aware Q-Former & Instruction-conditioned SFT & Task-specific visual feature extraction \\
Kosmos-2 & Interleaved tokens & SFT with grounding data & Bounding box grounding in text \\
LLaVA-RLHF & Linear projection & SFT + factuality-focused RLHF & First multimodal preference optimization \\
RLHF-V / mDPO & Various & DPO on hallucination pairs & Fine-grained hallucination reduction \\
LLaVA-OneVision & AnyRes encoding & Multi-stage SFT & Unified image / video / multi-image \\
\bottomrule
\end{tabularx}
\end{table}

The evolutionary trajectory of Thread~5 demonstrates the capability progression of VLMs from ``seeing'' to ``understanding'' to ``reasoning'':

\begin{enumerate}[leftmargin=2em]
  \item \textbf{Visual Instruction Following} (2022--2023):
    Flamingo established the paradigm of visual in-context learning,
    LLaVA and InstructBLIP demonstrated that visual instruction tuning
    can unlock multimodal capabilities at extremely low cost,
    igniting the research wave of VLM post-training.

  \item \textbf{Alignment and Hallucination Mitigation} (2023--2024):
    LLaVA-RLHF and RLHF-V introduced preference optimization into the multimodal domain,
    benchmarks such as POPE and HallusionBench exposed the systematic nature of hallucination,
    and methods such as RLAIF-V and mDPO provided scalable alignment solutions.

  \item \textbf{Unification and Reasoning} (2024--2025):
    LLaVA-OneVision achieved unified processing of images, videos, and multi-image inputs,
    multimodal reasoning extended Thread~4's chain-of-thought capabilities to the visual domain,
    and omni models began pursuing cross-modal capability consistency.
\end{enumerate}

Multimodal post-training endows models with the ability to ``see'' and ``understand'' the visual world,
but the true value of this capability lies in laying the foundation for \textbf{action}.
When VLMs can understand screenshots, localize UI elements, and interpret charts,
they are just one step away from ``autonomously operating in real environments.''

\paragraph{Preview of the Next Chapter.}
When VLMs can understand screenshots, localize UI elements, and interpret charts,
they possess the perceptual prerequisites for autonomous operation in real environments.
Thread~6 will trace the evolutionary path of agentic capabilities ---
from RL-based tool use to coding agents to computer use,
demonstrating how models transform from ``passive text generators'' into ``active world interactors.''
